\chapter{Introduction}
\label{ch:introduction}

Research groups commonly expose their work through publication lists. Such lists answer what has been published, but leave substantial interpretation work to the reader: locating a paper, obtaining its full text, understanding its contribution, relating it to other work, and deciding whether it can support a thesis or collaboration. The precursor TTLAB system automated collection, display, summarisation, notifications, and podcast preparation over publication metadata and abstracts, while identifying retrieval-augmented querying as future work \cite{narine2025automating}. The present project extends that direction into a full-text research-intelligence prototype.

The broader scholarly-discovery landscape includes autonomous digital libraries such as CiteSeer \cite{giles1998citeseer}, researcher-centred services such as ArnetMiner \cite{tang2008arnetminer}, large literature graphs such as S2ORC \cite{ammar2018literaturegraph}, and open bibliographic infrastructure such as OpenAlex \cite{priem2022openalex}. These systems demonstrate the value of linked metadata and retrieval at scale. This project addresses a smaller and different setting: a curated laboratory corpus, local operation, explicit page/chunk evidence, student-facing extension advice, and review before public exposure.

The implemented application is not treated as proof that all of those goals have been achieved. Instead, this thesis adopts a strict evidential boundary. A capability is called implemented only when supported by code, persisted state, an executed test or probe, or a captured interface. A quality claim requires appropriate evaluation evidence. Proposed experiments and missing institutional decisions are marked explicitly.

\section{Aim}

The aim is to design, implement, and reproducibly characterize a source-traceable research-intelligence prototype for TTLAB publications, while identifying the evidence still required to establish retrieval effectiveness, answer faithfulness, recommendation usefulness, and readiness for public deployment.

\section{Objectives}

The project objectives are to:

\begin{enumerate}[label=O\arabic*.,leftmargin=13mm]
  \item construct a deterministic pipeline from publication records and local PDFs to page-aware, searchable chunks;
  \item provide searchable publication, topic, and author views over a curated corpus;
  \item implement retrieval and citation-returning question answering with visible evidence and grounding warnings;
  \item produce source-linked thesis-extension recommendations that distinguish paper evidence from system suggestions;
  \item expose generated outputs and metadata to administrative review rather than silently presenting them as verified; and
  \item define and support a reproducible evaluation method that separates implementation checks from empirical effectiveness.
\end{enumerate}

\section{Contributions}

The defensible contributions of the audited snapshot are:

\begin{itemize}
  \item a modular local software artefact joining PDF extraction, page-aware chunking, hybrid retrieval, source-linked generation, exploration, and review;
  \item a structured recommendation contract separating baseline work, evidence status, inferred or explicit gaps, generated extension scope, risks, skills, and evaluation suggestions;
  \item an auditable evidence package containing machine-readable repository and runtime snapshots, reproducible figure sources, build commands, and an author-review checklist; and
  \item a candid characterization of coverage, data quality, index consistency, and evaluation gaps in a real laboratory corpus.
\end{itemize}

No algorithmic novelty, retrieval superiority, factual-accuracy rate, recommendation benefit, or deployed production service is claimed.

\section{Thesis structure}

\Cref{ch:problem} defines the problem and scope. \Cref{ch:rqs} presents the research questions and traceability plan. \Cref{ch:literature} reviews relevant scholarship. \Cref{ch:methodology,ch:architecture,ch:data,ch:implementation} describe the evidence method, design, corpus, and implementation. \Cref{ch:experiments,ch:results} separate planned evaluation from currently observed results. \Cref{ch:discussion,ch:limitations,ch:conclusion} interpret the evidence, state threats and ethical requirements, and conclude.
